\clearpage\thispagestyle{empty}\begingroup\small\raggedright
\noindent{\Large\bfseries Translation sources and terminology\par}\medskip
Historical text: Carl Friedrich Gauss, \emph{Werke}, Volume III, Göttingen, 1866. These English and German texts are AI-translated from the original-language edition. The following human translations are consulted for the stated terminology and usage; their running text is not reproduced. Gauss’s original mathematical formulations remain controlling.
\par\medskip\textbf{Paul Taylor; Bernard Leak.} Taylor’s English translation of Gauss’s second proof of the fundamental theorem of algebra, with corrections by Leak, 1983. Related English polynomial, divisor, degree and proof terminology; the 1815 paper is not a parallel translation of the 1799 dissertation. Its later use of complex quantities does not replace Gauss’s definition of imaginary quantities here.\par
\url{https://www.paultaylor.eu/misc/gauss.html}
\par\medskip\textbf{H. Maser.} German translator and editor, \emph{Carl Friedrich Gauss’ Untersuchungen über höhere Arithmetik}, Berlin: Julius Springer, 1889. Related German terminology for algebraic functions, roots, divisibility, quotients and degrees; not a parallel translation of this dissertation.\par
\url{https://archive.org/details/carlfriedrichga00gausgoog}
\par\medskip\textbf{James Caddall Morehead and Adam Miller Hiltebeitel.} Translators of Gauss’s \emph{General Investigations of Curved Surfaces of 1827 and 1825}, Princeton University Library, 1902; Project Gutenberg 36856, TeX edition. Related English geometrical, coordinate and infinitely small quantity terminology; not a parallel translation of the 1799 dissertation.\par
\url{https://www.gutenberg.org/ebooks/36856}
\par\medskip\textbf{AI authorship and publication rights.} AI-translated using OpenAI GPT-6.1 Sol, with AI-produced typesetting. Gauss is the author of the historical mathematical text. The named human translations are credited for their stated terminology and usage roles.
\par The editions are available free of charge for non-commercial, purely educational purposes. Publication rests on the publisher’s presumption that the respective rights holders would not object; this is not established permission. Free access to a source does not itself authorize derivative publication. Rights in the originals and any protected derivative contributions remain with their respective rights holders. No project-wide licence or public-domain dedication is granted here. Rights holders may raise objections through the repository’s issue page. The publisher undertakes to take down the affected material upon a rights holder’s request.\par
\url{https://github.com/KokunoYumeto/modern-latex-manuscripts/issues}
\par\medskip Scholarly citation metadata: \texttt{CITATIONS.bib}.
\par\endgroup\clearpage
